\section{Pre-registration and deviations}\label{app:prereg}

The pre-registration fixes every definition, constant, split, hypothesis and selection rule. After it was
written, changes could only be appended to its deviations log, each with its date, the section it touches, the
reason and whether it came before or after the freeze; the freeze record holds the log's length at the freeze, and
the lock refused to release the test outcomes unless the log was an unchanged prefix of the current one.
Table~\ref{tab:a-deviations} gives every entry in one line. Most entries make precise what the pre-registration
left open (how a filter reaches its budget, which events a sample holds, how a variable is scaled) and were written
before the freeze. The material deviations are these.
\begin{itemize}[leftmargin=*, itemsep=1pt]
\item \emph{Amendments before any data.} Convergence to the news-implied value, a wider question battery and the
  firm-level explanation analysis were added on the day the pre-registration was written.
\item \emph{Masking.} The masking rules were widened several times before the freeze, after independent checks
  found residual names, dates and price moves; non-English articles, move explainers and running market stories
  were excluded; residues below half a percent of states were accepted and disclosed.
\item \emph{The commercial models} failed the pre-registered screens. Instead of dropping them, the three with
  the highest second-stage IC were kept as screen-failing, model-specific comparators (before the freeze).
\item \emph{Training.} The trained instruments are one seed and one epoch each, not the pre-registered three-seed
  ensemble, because of a week's free GPU quota; the other seeds are robustness runs, fixed before unblinding and
  not yet run.
\item \emph{H1b} (the IC with the move after the open) was added to the primary family before the freeze, because
  the strategy that enters at the open needs it.
\item \emph{Company identification} was re-run with activity-matched distractors before the freeze, because a
  guess from news activity alone beat chance in the original design.
\item \emph{Items not run.} Three robustness variants needed new model calls and were not run (a summary-text
  variant, vendor tags in the text and a 09:30 news cutoff); the 08:00 cutoff is approximated from time stamps.
  Of the ablation requests, \res{ablation-failed} were rate-limited and not retried.
\item \emph{The canary.} After the freeze, a stability check that repeats earlier requests fired for span-02
  before production. Its build string was unchanged and the drift was close to its own repeat noise, so production
  went ahead under a rule recorded before the end-of-production check: a model counts as updated only if its build
  changes or two independent checks fire. At the end of production neither condition held for any model.
\item \emph{The price-only baseline's post-close predictions} were moved after unblinding, because one of their
  inputs, the day's overnight move, is a test-period outcome; the model was unchanged and first reproduced its
  validation predictions exactly.
\item \emph{Analysis and strategies.} The operationalization of the frozen specification and of the strategies
  was written before the outcomes were read; three code fixes found when the code first met the test data changed
  no pre-registered number.
\item \emph{Exploratory work after unblinding} is dated in the log. The conditioning study's validation selection
  was committed before its single test run; its two definitional choices (the VIX close of the previous session,
  read from a registered copy of the market data, and payroll release dates added from the Bureau of Labor
  Statistics) were made on validation data.
\item \emph{The forward test} was pre-registered after the paper's analyses and before any forward-period data
  existed (Section~\ref{sec:limits}).
\end{itemize}

\input{\tabdir/t34_prereg_deviations.tex}

\section{Notation}\label{app:notation}

{\small
\begin{longtable}{@{}>{\raggedright\arraybackslash}p{1.35in}>{\raggedright\arraybackslash}p{4.95in}@{}}
\caption{Notation.}\label{tab:notation} \\
\toprule
Symbol & Meaning \\
\midrule
\endfirsthead
\multicolumn{2}{@{}l}{\textit{Table \ref{tab:notation}, continued}} \\
\toprule
Symbol & Meaning \\
\midrule
\endhead
\bottomrule
\endlastfoot
D0, D1, D2 & Decision points: 09:15~ET before the open (pre-open reading); 10:00~ET after the first 30-minute bar
and after the close (post-reaction readings). \\
val-A, val-B, test & 2025-07-01 to 2025-09-30; 2025-10-01 to 2025-12-31; 2026-01-02 to 2026-07-31. \\
$r^{cc}, r^{on}, r^{id}$ & Close-to-close, overnight (close to open) and intraday (open to close) returns, with
$1 + r^{cc} = (1 + r^{on})(1 + r^{id})$; Equation~\eqref{eq:split}. \\
$\AR^{k}$, $\sigma^{k}_{i,t-1}$ & Abnormal return (market and leave-one-out sector removed, betas over 252
sessions to $t-1$); its ex-ante scale (1.4826 times the median absolute deviation over 126 sessions to $t-1$,
floored at the day's cross-sectional 10th percentile). \\
$z^{k}$; $z_{cc}, z_{on}, z_{id}$ & $\AR^{k}/\sigma^{k}_{i,t-1}$, the move in typical daily moves (clipped at
$\pm 10$). \\
$y_0, y_1, y_2$ & Outcomes after D0, D1 and D2: $z_{id}$; the move from 10:00 to the close; the move from the next
open to the close of $t+5$, each over its ex-ante volatility. \\
$e_5$ & The move from the previous close to the close of $t+4$ over $\sigma^{cc}\sqrt{5}$: the training label and
the calibration target. \\
$G$, $\Glvl$ & The observed reaction at D1 or D2 in typical-move units; its level version, divided by
$\sigma^{cc}$ without a square-root-of-time factor. \\
$\Nhat$, $\Ntilde$ & The news-implied move $\sum_k p_k m_k$ over the nine levels of the pre-open reading; its
calibrated version (an isotonic map to $e_5$ fitted on val-A), which ranks the pre-open strategy. \\
$V$, $U$ & The news-implied value (an isotonic map of $\Nhat$ to the level move to $t+4$, fitted on val-A); the
remaining gap $V - \Glvl$ at D1 or D2. \\
$J$ & The post-reaction judgment $P(\text{much too small}) + P(\text{too small}) - P(\text{too large}) -
P(\text{much too large})$. \\
Under, Fair, Over & Post-reaction classes from $(\Nhat - G)\,\mathrm{sign}(\Nhat)$ for $|\Nhat| \ge 0.5$: Under
above 0.5, Over below $-0.5$, Fair between. \\
$\CARlvl_{h}$ & The level cumulative abnormal return from the previous close to rung $h$. \\
$\kappa_h$, $\theta_V$ & The convergence slope at rung $h$ (Equation~\eqref{eq:ladder}); the reliability of $V$
(Equation~\eqref{eq:attenuation}). \\
$\lambda_h$, $\psi_h$, $h^{*}$ & The share of the gap closed by rung $h$, the convergence to a shrunken value and
the half-life (Equation~\eqref{eq:gap}). \\
$\IC_t$, $\overline{\IC}$ & The daily Spearman information coefficient and its mean over sessions
(Equation~\eqref{eq:ic}). \\
$R^{2}_{\mathrm{OOS}}$, $\Delta R^{2}$ & The out-of-sample $R^2$ and a step's increment
(Equation~\eqref{eq:r2}). \\
$\mathrm{Acc}_m$, $p_0$ & Company-identification accuracy and the matched baseline
(Equation~\eqref{eq:memorization}). \\
ECE & Expected calibration error of the top class of three (a fall, no move or a rise, split at $\pm 0.5$
typical moves against $e_5$), 15 equal-width bins. \\
DSR, $\mathrm{SR}_0$ & The deflated-Sharpe probability and its benchmark (Equation~\eqref{eq:dsr}). \\
H1, H1b, H2, H3, H4 & Primary hypotheses (Holm at 0.05). \\
H5, H6, H7, H2-C & Secondary hypotheses (Benjamini--Hochberg at $q = 0.10$): explaining moves, style, the
instruments, convergence. \\
F1 to F7 & The forward test's hypotheses (Section~\ref{sec:limits}). \\
\end{longtable}
}

\section{The deflated Sharpe ratio}\label{app:dsr}

For a book with $T$ daily net returns, a per-period Sharpe ratio $\widehat{\mathrm{SR}}$, skewness
$\hat\gamma_3$ and kurtosis $\hat\gamma_4$, the deflated-Sharpe probability \citep{bailey2014} is computed as
\begin{equation}\label{eq:dsr}
\mathrm{DSR} = \Phi\!\left(\frac{(\widehat{\mathrm{SR}} - \mathrm{SR}_0)\sqrt{T-1}}
{\sqrt{1 - \hat\gamma_3\,\widehat{\mathrm{SR}} + \tfrac{\hat\gamma_4 - 1}{4}\,\widehat{\mathrm{SR}}^{2}}}\right),
\qquad
\mathrm{SR}_0 = \sqrt{\mathcal{V}}\left[(1 - \gamma)\,\Phi^{-1}\!\left(1 - \frac{1}{N}\right)
+ \gamma\,\Phi^{-1}\!\left(1 - \frac{1}{N\mathrm{e}}\right)\right],
\end{equation}
where $\Phi$ is the standard normal distribution function, $\gamma$ the Euler--Mascheroni constant, $N$ the
number of trials and $\mathcal{V} = (1 + \widehat{\mathrm{SR}}^{2}/2)/(T-1)$ the variance of a Sharpe-ratio
estimate \citep{lo2002}; the analytic variance is used because most of the ledger's other trials are information
coefficients, recalls and losses, not Sharpe ratios. $\mathrm{SR}_0$ is the Sharpe ratio that the best of $N$
independent trials without skill would reach in expectation, so it rises with $N$: a book must beat it, not zero.
With the \res{ntrials-strat} trials of the ledger when the pre-registered strategies were evaluated and
\res{sconv-ndays} test days, $\mathrm{SR}_0$ is \res{dsr-sr0} per day for the convergence strategy, far above any
book's daily Sharpe ratio. The exploratory books of Section~\ref{sec:exploratory} were deflated with all
\res{ntrials-total} trials of the ledger after the exploratory study, which raises the bar further.

\section{Additional pre-registered results}\label{app:results}

This appendix gives the pre-registered results that the main text cites but does not show: every model's IC
and cost (Table~\ref{tab:a-model-comparison}), the robustness variants (Figure~\ref{fig:a-robustness}), the
heterogeneity cuts (Table~\ref{tab:a-heterogeneity}), the secondary hypotheses (Table~\ref{tab:a-secondary}) and
calibration (Table~\ref{tab:a-calibration} and Figure~\ref{fig:a-reliability}). Further tables and figures are
in the replication package.

\input{\tabdir/t14_model_comparison.tex}

\begin{figure}[htbp]
\centering
\includegraphics{\figdir/figA5_robustness.pdf}
\caption{\textbf{No pre-registered robustness variant changes a conclusion: Laya-Fin's H1 and H1b stay within
their intervals of zero, and d1-ultra's H1 stays positive in every variant, from \res{d1-h1-pre04} when the
latest article precedes 04:00 to \res{d1-h1-0815} when it arrives between 08:00 and 09:15.} Test period, D0
events. Each variant raw (filled) and with placebo-adjusted targets (hollow), below the main estimate (diamond,
from the stationary block bootstrap), with 95\% intervals ($\pm 1.96$ Newey--West standard errors for the
variants). (a) and (d), both H1, share a scale; (b) and (c) have their own. d1-ultra is
model-specific.}\label{fig:a-robustness}
\end{figure}

\input{\tabdir/t5_heterogeneity.tex}

\input{\tabdir/t4_secondary.tex}

\input{\tabdir/t15_calibration.tex}

\begin{figure}[htbp]
\centering
\includegraphics{\figdir/figA4_reliability.pdf}
\caption{\textbf{Laya-Fin's three-class forecasts are calibrated but compressed (expected calibration error
\res{lf-ece}); the commercial models are overconfident (\res{drex-ece} to \res{span-ece}).} Reliability of the
pre-open reading's three classes (a fall, no move or a rise, split at $\pm 0.5$ typical moves) against the
five-day move $e_5$ on the test period: the hit rate of the top class by bin of its stated probability (15 bins;
the dot area grows with the number of events in the bin; the diagonal is perfect calibration). The anchor's
near-constant readings sit in two bins. $^{\ddagger}$span-02: see Figure~\ref{fig:validation-test}. The commercial
models are model-specific.}\label{fig:a-reliability}
\end{figure}

\FloatBarrier
\section{Exploratory: the conditioning study}\label{app:exploratory}

Figure~\ref{fig:a-conditioning} shows every (model, condition) pair of the post hoc conditioning study of
Section~\ref{sec:exploratory}, with the \res{cond-selected} pairs that the validation rule selected shaded. It is exploratory:
the study was requested after unblinding, and the figure supports no claim.

\begin{figure}[htbp]
\centering
\includegraphics{\figdir/figA18_conditioning_map.pdf}
\caption{\textbf{Exploratory (post hoc): none of the \res{cond-selected} conditions selected on two validation
quarters held on the test; \res{cond-reversed} changed sign.} For each (model, condition) pair, the difference
inside minus outside the condition of the mean $\mathrm{sign}(\Nhat)\cdot z_{id}$ (typical daily moves) on val-A
(circles), val-B (squares) and the test (dots), with 95\% intervals; shaded rows were selected by the rule fixed
before the test run (positive with $t \ge 1$ in both validation quarters and a pooled $t \ge 2$). $^{\dagger}$Vendor
fields are retrospective labels and only split events descriptively. ``(seen)'' marks conditions whose test values
the authors' assistant had seen before the study, which were never selected. Requested after unblinding; no
claim.}\label{fig:a-conditioning}
\end{figure}

\section{The replication package}\label{app:replication}

The replication package is the project's repository \placeholder{repository URL to be provided by the authors}.
It holds the pre-registration with its amendments and its deviations log; the freeze record (the hash, its time
and the version-control tag); the code of every stage, from the price panel and the masked news states to the
model calls, the analysis and the strategies, with its tests; the committed outputs and a report of every stage;
the append-only, hash-chained ledgers of API spend, trials, test-data access and data requests; and the sources of
this paper. Further tables and figures are in the replication package. One command rebuilds every table, figure
and caption of the paper from the committed outputs (plus three stored strategy files that the strategy stage
regenerates) and checks each against a committed digest; another writes the macros through which every result
number in this paper enters the text, each a cell of those tables, and checks that no result number in the
sources is typed by hand. Neither needs raw data, an API, a GPU or the lock. The news text is licensed and is not
redistributed; \placeholder{data availability statement to be provided by the authors}.
